
\documentclass[fleqn,10pt]{olplainarticle}
% Use option lineno for line numbers 

\usepackage{algorithm}
\usepackage{algorithmic}
\begin{document}

\subsection{System Model and Assumptions}
\label{subsec:system-model}

We consider a single intersection instrumented with connected and autonomous vehicles (CAVs) that cooperate to agree on a safe and fair crossing order. Our protocol runs in three modes---BFT consensus, abortable small-$n$ consensus, and a time-slotted fallback---all under the following common model.

\paragraph{Network and timing.}
We assume a \emph{partially synchronous} communication model: there exists an unknown global stabilization time after which all messages between vehicles within communication range are delivered within a bounded delay $\Delta$. Each vehicle $v$ has access to a local hardware clock $\mathsf{clk}_v(t)$, and clocks are synchronized using a trusted time-synchronization service (e.g., GNSS + infrastructure support). We assume a known bound $\epsilon$ on clock skew:
\[
  \forall u,v:\;\bigl|\mathsf{clk}_u(t) - \mathsf{clk}_v(t)\bigr| \le \epsilon.
\]
In fallback mode, time is partitioned into \emph{frames} of length $F$, and each frame is further divided into \emph{slots} of length $\delta$. We choose $F$ and $\delta$ such that
\[
  \delta \gg \epsilon + \Delta + \text{(braking / sensing margin)},
\]
so that slots assigned to conflicting maneuvers do not overlap for honest vehicles that respect the slotting rule.

\paragraph{Vehicles, identities, and fault model.}
Each physical vehicle $v$ is provisioned with a globally unique, certification-authority (CA) issued public key $\mathsf{pk}_v$ and corresponding private key $\mathsf{sk}_v$. Certificates bind $\mathsf{pk}_v$ to a physical vehicle identity and prevent Sybil attacks (i.e., a vehicle cannot cheaply create multiple identities). All vehicles can verify signatures and certificates of other vehicles.

We assume at most one Byzantine vehicle among the set of vehicles that are simultaneously within communication range of the intersection:
\[
  f = 1.
\]
A Byzantine vehicle may deviate arbitrarily from the protocol: send conflicting messages, omit messages, or refuse to sign. However, the adversary cannot break standard cryptographic primitives (e.g., cannot forge signatures) and cannot forge CA-issued certificates. Physical aspects such as a vehicle deliberately driving through a red phase or permanently blocking the intersection are handled by separate safety and enforcement mechanisms and are out of scope for the protocol analysis.

\paragraph{Intersection identity and localization.}
Each intersection $I$ is associated with:
\begin{itemize}
  \item a unique logical identifier $\mathsf{id}_I$ (e.g., derived from map coordinates), and
  \item a geographic region $R_I$ that captures the spatial extent of the intersection.
\end{itemize}
Vehicles use on-board localization (GNSS, maps) to determine when they are within $R_I$. We assume that a vehicle cannot arbitrarily claim to be in a different intersection region without being detected by the physical safety system; i.e., localization errors are bounded and cannot be exploited to move a certificate from one intersection to another. All quorum certificates described below are bound to a specific intersection identifier $\mathsf{id}_I$.

\paragraph{Quorum certificates (QCs).}
A \emph{quorum certificate} (QC) for intersection $I$ and epoch $e$ has the form
\[
  \mathsf{QC} = \Bigl(\mathsf{id}_I,\; e,\; \mathsf{order},\; \Sigma\Bigr)
\]
where:
\begin{itemize}
  \item $\mathsf{order}$ is a finite sequence of vehicle identifiers describing a crossing order, and
  \item $\Sigma$ is a set of signatures over $(\mathsf{id}_I, e, \mathsf{order})$.
\end{itemize}
We distinguish the following cases:
\begin{itemize}
  \item For $n \ge 4$ participants (the BFT mode with $f=1$), a QC is valid if $|\Sigma| \ge 3$, i.e., if it contains signatures from at least three distinct vehicles present at the intersection.
  \item For $n = 3$ participants, we require \emph{unanimity}: a QC is valid only if all three vehicles sign the same $\mathsf{order}$, i.e., $|\Sigma| = 3$. If any vehicle withholds a signature or sends an inconsistent proposal, no QC is formed and the small-$n$ consensus attempt aborts.
  \item For $n = 2$ participants, the protocol does not form a full QC but instead uses bilateral \emph{go/yield} certificates as described below.
\end{itemize}
In all cases, vehicles verify that $\mathsf{QC}$ is correctly signed, bound to $\mathsf{id}_I$, and that $\mathsf{order}$ satisfies the protocol's fairness and safety constraints before accepting it.

\paragraph{Waiting certificates (waitingQCs).}
When consensus aborts, or when a vehicle remains waiting while others are allowed to cross, we issue weaker certificates of the form
\[
  \mathsf{wQC} = \Bigl(\mathsf{id}_I,\; e,\; v,\; \mathsf{meta},\; \Sigma'\Bigr),
\]
where $v$ is the vehicle identifier, $\mathsf{meta}$ may include local arrival time and approach lane, and $\Sigma'$ is a set of signatures attesting that ``vehicle $v$ was present at intersection $I$ in epoch $e$ and has not yet crossed.'' We require $|\Sigma'| \ge 2$ whenever possible (e.g., in $n=2$ and $n=3$ scenarios), so that no single vehicle can unilaterally mint a waitingQC for itself. In the $n=2$ case, both vehicles sign each other's waitingQCs when one yields and remains behind.

WaitingQCs do not directly encode a complete crossing order, but they provide cryptographic evidence of \emph{waiting history} that influences priority in future decisions and in the fallback time-slotted regime.

\paragraph{Epochs and certificate validity.}
Time at an intersection is divided into \emph{epochs} $e = 0,1,2,\dots$. Each QC and waitingQC is tagged with the epoch in which it was created. Certificates are only valid for a bounded time window $T_{\mathit{expire}}$ (e.g., on the order of a few minutes):
\[
  \text{A QC or waitingQC is valid at time } t \text{ only if } t - t_{\mathit{issue}} \le T_{\mathit{expire}}.
\]
An epoch is reset (and all prior QCs and waitingQCs are ignored for fairness) when:
\begin{itemize}
  \item the intersection has been idle (no vehicles within $R_I$) for longer than a configured threshold, or
  \item the latest QC and associated waitingQCs have all expired.
\end{itemize}
This prevents vehicles from replaying stale certificates obtained at the same intersection long in the past to cut ahead of a new queue.

\paragraph{Modes of operation.}
Under these assumptions, the protocol operates in three modes:
\begin{itemize}
  \item \emph{BFT mode} (typically $n \ge 4$): vehicles run a Byzantine fault-tolerant consensus protocol that produces QCs satisfying safety and a strong fairness property.
  \item \emph{Small-$n$ abortable mode} ($n \in \{2,3\}$): vehicles attempt to agree unanimously on a crossing decision; on success, they issue a QC or bilateral go/yield certificates; on failure, no vehicle moves and waitingQCs are issued instead.
  \item \emph{Fallback time-slotted mode}: if consensus fails to make progress, vehicles use their synchronized clocks, QCs, and waitingQCs to compute priority-aware time slots for crossing, while relying on onboard sensing to ensure physical safety.
\end{itemize}
In all three modes, we assume that at most one vehicle within communication range is Byzantine and that honest vehicles follow the protocol and do not intentionally violate physical safety constraints.

\end{document}